\documentclass[10pt,a4paper]{amsart}
\usepackage[margin=19mm]{geometry}
\usepackage{amsmath,amssymb,mathtools}
\usepackage[scr=rsfs,cal=euler]{mathalfa}
\usepackage{xfrac}
\usepackage[unicode,hidelinks,bookmarks=false]{hyperref}
\newcommand{\ie}{i.e.\ }
\newcommand{\cf}{cf.\ }
\newcommand{\resp}{respectively\ }
\newcommand{\Subsection}{\subsection}
\providecommand{\nobreakdash}{\nobreak\hbox{-}}
\setlength{\emergencystretch}{2em}
\multlinegap0pt
\usepackage{tikz}
\usepackage{amssymb}
\usepackage{mathtools}
\usepackage[shortlabels]{enumitem}
\usepackage[normalem]{ulem}
\usepackage{xfrac}
\usepackage{xcolor}
\newtheorem{proposition}{Proposition}
\newcommand{\PP}{ \ensuremath{\mathbb{P}}}
\title{Weak and Strong Lefschetz Properties for Hartshorne-Rao Modules of Curves in $\mathbb P^3$}
\author{Juan Migliore, Uwe Nagel, Chris Peterson, Ettore Teixeira Turatti}
\date{Source: arXiv:2605.19434v1 (CC BY 4.0)}
\begin{document}
\maketitle

\noindent\emph{Source and licence.} Weak and Strong Lefschetz Properties for Hartshorne-Rao Modules of Curves in $\mathbb P^3$. Version arXiv:2605.19434v1 (CC BY 4.0), distributed by arXiv under the Creative Commons Attribution 4.0 International licence, http://creativecommons.org/licenses/by/4.0/, as recorded in the arXiv licence metadata for this exact version and shown on the abstract page https://arxiv.org/abs/2605.19434v1. Full licence notice: \texttt{licenses/arxiv-2605.19434-CC-BY-4.0.txt}.\par\smallskip

\noindent\textit{Editorial scope.} Counted unit: the article's proposition Z2 (source0.tex lines 615--621). The article's complete original proof of it is delivered with it (source0.tex lines 623--830, one \texttt{proof} environment of 208 lines). No mathematics is written, repaired or re-derived: the statement and the proof are the article's own lines, in the article's own notation, and no retained line is substituted or re-flowed.  Retained uncounted as the source-local context the proof needs: lines 474--482.  Nothing was omitted from any retained span, so the package carries no omission record.  No retained line contains a citation, so the wrapper carries no bibliography and no bibliography entry.  No retained line contains a figure environment, an image-inclusion command or drawing code, so no image is delivered and the package declares no asset dependency.  Editorial formatting only, in addition: an AMS \texttt{amsart} preamble that restates the article's own declarations and macros used by the retained lines, namely the environments \texttt{proposition} on separate counters, together with the standard packages those lines need.  The retained environments print in this document as Proposition 1 in order of appearance, whereas the article prints the same items with the numbering of its own full document.  The complete native source of the article as distributed in the arXiv e-print for this exact version is bundled unchanged as \texttt{sources/200914/source0.tex}.
\begin{proposition} \label{GenFlatFat}
    Let $Z_{s,m}$ be a general union of $s$ flat fat points of multiplicity $m$ in $\mathbb P^2$. Assume that $1 \leq m \leq 3$. If $m = 1$ or $m=2$, assume $s \geq 1$, and if $m = 3$ assume that $s \geq 3$. 
    
    Then $Z_{s,m}$ has generic Hilbert function, i.e.
    \[    
    h_{Z_{s,m}}(j) = \min \left \{ ms, \binom{j+2}{2} \right \}.
    \]

\end{proposition}

\begin{proposition} \label{Z2}
Let $C \subset \PP^3$ be a set of $r$ general (disjoint) lines and let $H$ be a general hyperplane, defined by a linear form $L$. Let $Z_2$ be the zero-dimensional scheme defined by the saturation of the ideal $I_C + (L^2)$. Then $Z_2$ is in generic position on the surface $H^2$ defined by $L^2$, {i.e., 
\[
h_{Z_2} (j) = \min \{ 2r, (j+1)^2\}. 
\]
}
\end{proposition}

\begin{proof} 
Note that $Z_2$ has degree $2r$ by B\'ezout's theorem. We define the scheme $Z_1$ by
\[
I_{Z_1}  =  I_{Z_2} : (L) \subset R.
\]
Note that we also have 
\[
I_{Z_1 |H}  =  (I_{Z_2} + (L))^{sat}/(L) \subset R/(L) ,
\]
\[
 [I_{Z_2} + (L)]/(L) \cong I_{Z_2} / [I_{Z_2} \cap (L)] \cong  I_{Z_2} / [L \cdot [I_{Z_2} : (L) ]]
\]
and we have the exact sequence
\[
0 \rightarrow [I_{Z_2} : (L)](-1) \stackrel{\times L}{\longrightarrow} I_{Z_2} \rightarrow [I_{Z_2} + (L)]/(L) \rightarrow 0
\]
or (by sheafifying and taking cohomology)
\begin{equation} \label{usual exact seq}
0 \rightarrow I_{Z_1}(-1) \stackrel{\times L}{\longrightarrow} I_{Z_2} \rightarrow I_{Z_1 |H} \rightarrow H^1_*(\mathcal I_{Z_1})(-1) \rightarrow H^1_* (\mathcal I_{Z_2}) \rightarrow H^1_* (\mathcal I_{Z_1}) \rightarrow 0
\end{equation}
where $H^1_*$ refers to a direct sum over all twists. Here we have used the fact that if $X$ is a zero-dimensional scheme on a surface $F$ of degree $d$  in $\PP^3$ then we have the short exact sequence
\[
0 \rightarrow R(-d) \rightarrow I_X \rightarrow I_{X|F} \rightarrow 0
\]
so sheafifying and taking cohomology we have $H^1_*(\mathcal I_X) \cong H^1_*(\mathcal I_{X|F})$.



We argue by specialization. First, we give the general framework, then we give the specific details for our result.

1. Assume without loss of generality that $H$ is defined by $x_3$. Consider a component of $Z_2$ supported at a point $P$. This  component is defined by $I = (\ell_1, \ell_2, x_3^2)$. Then $I$ is determined by (i) the ideal of the point $P$, that is, $I_P = (\ell_1, \ell_2, x_3)$, and (ii) the line defined by $[I]_1$, the “tangent line,” because if $m$ is any linear form such that $(\ell_1, \ell_2, m) = (\ell_1, \ell_2, x_3)$ then 
$I = (\ell_1, \ell_2, m^2)$. 

2. Now we fix the reduced set $Z_1$ and we consider sets of $r$ lines, where for each point $P$ of $Z_1$ a line in $\mathbb P^3$ is chosen through $P$. Each such choice determines an ideal with local components as in 1. 
 Since $Z_2$ is determined by this information, it is specified by $Z_1$ plus a line through each point of $Z_1$.

We can take as our parameter space, then, the variety
\[
\underbrace{(\mathbb P^2) \times \dots \times  (\mathbb P^2)}_{\hbox{$r$ copies}}.
\]
If the line {$\lambda$ } through $P$ is general (we use $L$ to denote the linear form defining $H$), then the component of $Z_2$ supported at $P$ is defined by $I_\lambda + x_3^2$.
Thus, our $Z_2$ is a general set in the parameter space. 
However, if the line $\lambda$ is contained in $H$, i.e. $I_\lambda = (m_1, x_3)$, then the local component at $P$ is defined by 
$(m_1, x_3, m_2^2)$ with some $m_2$ chosen so that $I_P = (m_1, x_3, m_2)$. 

Our specialization, then, will come from moving general lines onto the plane $H$.






Let $t$ be the smallest integer such that $r \leq \binom{t+1}{2}$. Let $\alpha = \binom{t+1}{2} - r$. 
This means that the $h$-vector of $Z_1$ is
\[
(1,2,3,\dots, t-2, t-1, t-\alpha)
\]
where the $t-\alpha$ occurs in degree $t-1$. 
Note that since $Z_1$ can be taken to be a general set of $r$ points in $H$, $\dim [I_{Z_1 |H}]_{t-1} = \alpha$ and $t-1$ is the socle degree of the Artinian reduction of $\mathbb{K}[x,y,z]/I_{Z_1|H}$. If $r < \binom{t+1}{2}$, we also have that the initial degree of $I_{Z_1|H}$ is $t-1$; otherwise it is $t$.

For now assume that $r < \binom{t+1}{2}$.

Now we specialize. Choose any $\alpha$ of the nonreduced points of $Z_2$ and move the ``direction" as above so that it is a general direction on $H$. 


Let $Y_2$ be the resulting scheme. Note $Y_2$ still lies in the surface $H^2$. 
We now define two related schemes in $H$:

\begin{itemize}
    \item First, we define $X_1$ by 
    $I_{X_1} = ( I_{Y_2} + (L))^{sat}$. We also have
    \[
    I_{X_1|H} = [[I_{Y_2} + (L)]/(L)]^{sat}.
    \]
It defines the union of  $r-\alpha$ general points of $H$ and $\alpha$ general nonreduced, length 2 schemes on $H$, for a total degree of $r+\alpha = \binom{t+1}{2}$.
    By Lemma \ref{GenFlatFat}, generic multiplicity~$2$ flat fat points impose independent conditions in $H$.  Hence the $h$-vector of $X_1$ is
    \[
(1,2,3,\dots, t-2, t-1, t).
\] 


    \item Second, the scheme $X_2$ defined by $I_{X_2} = I_{Y_2} : (L)$ consists of (hence has the Hilbert function of) $r-\alpha$ general points in  $\PP^2 = H$.

\end{itemize}

\noindent The bottom line is that compared to the original scheme $Z_2$, here the zero-dimensional scheme on $H$ rises to degree $r+\alpha = \binom{t+1}{2}$ (the degree increases by $\alpha$), while the residual scheme drops in degree by $\alpha$. 

We will consider two cases, which we list first and then make the computations separately.

\vspace{.2in}

\noindent \underline{Case 1}: $\alpha \leq \frac{t}{2}$.

In this case the $h$-vectors of the ``residual" scheme defined by $ I_{X_2} = I_{Y_2} : (L)$, and that of the scheme defined by $I_{X_1} = [I_{Y_2} +(L)]^{sat}$ are
\[
(1,2,3,\dots, t-1,t-2\alpha) \ \ \hbox{ and } \ \ (1,2,3,\dots, t-1, t)
\]
where the last entries are both in degrees $t-1$.

\vspace{.2in}

\noindent \underline{Case 2}: $\alpha > \frac{t}{2}$.

In this case, for the residual scheme we have to reduce the total degree by $\alpha$ (compared to $Z_1$). Since $\alpha = (t-\alpha) + (2\alpha - t)$, we remove from the $h$-vector of $Z_1$ the $t-\alpha$ in degree $t-1$ plus an additional $(2\alpha-t)$ in degree $t-2$ to obtain the $h$-vector 
\[
(1,2,3, \dots, t-2, (t-1) - (2\alpha -t) ) = (1,2,3,\dots, t-2, 2t-2\alpha -1)
\]
for $X_2$, and we still have 
\[
(1,2,3,\dots, t)
\]
for $X_1$.
Here the last nonzero entries are in degrees $t-2$ and $t-1$ respectively. 

 For any integer $s$, we have the exact sequence
\begin{equation} \label{double2}
0 \rightarrow [I_{X_2}]_{s-1} \rightarrow [I_{Y_2 }]_s \rightarrow [I_{X_1 |H}]_s \rightarrow H^1(\mathcal I_{X_2}(s-1)) \rightarrow H^1 (\mathcal I_{Y_2}(s)) \rightarrow H^1 (\mathcal I_{X_1}(s)) \rightarrow 0.
\end{equation}

\noindent We observe that

\begin{itemize}

\item[(i)] $\dim [I_{X_1|H}]_s = 0$ for $s \leq t-1$.

\item[(ii)] $h^1(\mathcal I_{X_1}(s)) = 0$ for $s \geq t-1$.

\item[(iii)] In Case 1, $h^1(\mathcal I_{X_2} (s-1)) = 0$ for $s \geq t$.

\item[(iv)] In Case 2, $h^1(\mathcal I_{X_2} (s-1)) = 0$ for $s \geq t-1$.
\end{itemize}

We begin with Case 1. From exactness of (\ref{double2}) and (ii), (iii) we conclude

\begin{enumerate}

\item 
\[
\left.
\begin{array}{l}
h^1(\mathcal I_{X_2}(t-1)) = 0 \\
h^1(\mathcal I_{X_1}(t)) = h^1(\mathcal I_{X_1}(t-1)) = 0
\end{array}
\right \}
\Rightarrow
h^1(\mathcal I_{Y_2}(s)) = 0 \hbox{ for } s \geq t.
\]



\item $\dim [I_{Y_2}]_s = \dim [I_{X_2}]_{s-1}$ for $0 \leq s \leq t-1$, from (i).
    
\end{enumerate}

Since the $h$-vector of $X_2$ is
\[
(1,2,3,\dots,t-1,t-2\alpha)
\]
but now we consider $X_2$ as being in $\mathbb P^3$, we have for $s \leq t-1$ that 
\begin{equation} \label{dimY2}
\dim [I_{Y_2}]_s = \dim [I_{X_2}]_{s-1} =  \binom{s+1}{3} 
\end{equation}
(the ideal coincides with the ideal of the plane). 
Thus for $s \leq t-1$,
\[
h_{Y_2}(s) = \binom{s+3}{3} - \dim [I_{Y_2}]_s = (s+1)^2.
\]
This means that the $h$-vector of ${Y_2}$ has the form
\[
(1,3,5,7,\dots)
\]
until degree $t-1$. On the other hand, conclusion 1. above shows that the $h$-vector of $Y_2$ ends in degree at most $t$. Thus $Y_2$ has generic Hilbert function in the surface $H^2$. 


We now turn to Case 2, {where $\alpha > \frac{t}{2}$.} Now we have (again from the exactness of (\ref{double2}))

\begin{enumerate}

    \item 
    \[
    \left. 
\begin{array}{l}
h^1(\mathcal I_{X_2}(t-2)) = 0 \Rightarrow h^1(\mathcal I_{X_2}(t-1)) = 0 \\
h^1(\mathcal I_{X_1}(t)) = h^1(\mathcal I_{X_1}(t-1)) = 0
\end{array}
\right \}
\Rightarrow h^1(\mathcal I_{Y_2}(s)) = 0 \hbox{ for } s \geq t
    \]
    (using (ii) and (iv)).

    \item $\dim [I_{Y_2}]_s = \dim [I_{X_2}]_{s-1}$ for $0 \leq s \leq t-1$.
    
\end{enumerate}

This information is identical to that of Case 1, so the argument is also identical and (\ref{dimY2}) again holds.

Now, in both cases $Y_2$ is a specialization of $Z_2$, so for any $s$ we have by semicontinuity
\[
\dim [I_{Y_2}]_s \geq \dim [I_{Z_2}]_s .
\]
But both lie on the surface $H^2$, whose ideal has dimension $\binom{s+1}{3}$ in degree $s$. We have, using (\ref{dimY2}) for $s \leq t-1$, 
\[
\binom{s+1}{3} = \dim [I_{Y_2}]_s \geq \dim [I_{Z_2}]_s \geq \binom{s+1}{3}.
\]
Thus the original scheme $Z_2$ also has generic Hilbert function in $H^2$, since the $h$-vector ends in degree $t$.

Finally, when we introduced our specialization, we left open the case where $r = \binom{t+1}{2}$. In this case, no specialization is needed! $X_1$ already has $h$-vector $(1,2,\dots,t)$ and the argument goes as above.
\end{proof}

\end{document}
